\clearpage
\section{Tema de investigación}

El proyecto propone el diseño de un sistema híbrido para la predicción rápida de incendios forestales mediante la integración de simulación multiescala, procesamiento geoespacial, computación de alto rendimiento y recuperación de escenarios precalculados. La propuesta no busca desarrollar un nuevo simulador físico de incendios, sino aprovechar herramientas existentes y construir la capa de software necesaria para integrarlas en un flujo de trabajo orientado a la toma de decisiones durante una emergencia.

En la escala forestal se emplearía FARSITE, una herramienta capaz de estimar la propagación del fuego sobre terrenos extensos a partir de variables como el punto de ignición, la topografía, el tipo de combustible y las condiciones meteorológicas. Sus resultados permitirían identificar las zonas hacia las cuales probablemente avanzaría el incendio y determinar si este podría aproximarse a una interfaz urbano-forestal (WUI), donde la vegetación entra en contacto con viviendas, infraestructura y población. Para estas zonas críticas se utilizaría WFDS, que ofrece una simulación física más detallada de fenómenos como la dinámica del flujo, la combustión y la transferencia de calor, aunque con un costo computacional considerablemente mayor.

La combinación de ambas herramientas permitiría emplear un modelo rápido en las áreas donde no se requiere un nivel elevado de detalle y reservar el modelo de mayor costo para los sectores de mayor riesgo. Para resolver las limitaciones de tiempo durante una emergencia, el sistema contaría con una biblioteca de escenarios precalculados. Antes de que ocurra un incendio se ejecutarían numerosas simulaciones, variando el punto de inicio, la velocidad y dirección del viento, las condiciones ambientales, el terreno y las características del combustible. Los resultados de FARSITE y, en escenarios seleccionados de la WUI, de WFDS, quedarían almacenados para ser consultados posteriormente.

Cuando se presente un incendio, un usuario como un comandante de bomberos podría seleccionar su ubicación en un mapa e ingresar las condiciones actuales. El sistema compararía esos datos con los escenarios almacenados mediante un algoritmo de similitud que considere, entre otros factores, la distancia entre puntos de ignición, las diferencias en el viento, la humedad, el combustible y las características del terreno. De esta manera, recuperaría uno o varios escenarios cercanos y mostraría rápidamente una estimación visual de la propagación del fuego y de su posible comportamiento al alcanzar la zona urbana, sin esperar la ejecución de una simulación completa en tiempo real.

El aporte desde la Ingeniería de Sistemas se concentra en la construcción de un pipeline interoperable entre FARSITE y WFDS, incluyendo la transformación de coordenadas, geometrías, parámetros y formatos de entrada. También comprende el procesamiento de datos de Sistemas de Información Geográfica, como elevación, vegetación, topografía e infraestructura; la generación masiva y paralelizada de escenarios mediante computación de alto rendimiento; el almacenamiento y búsqueda eficiente de simulaciones; y el desarrollo de una aplicación de visualización accesible para el usuario final. Así, el proyecto integraría modelos existentes y convertiría sus resultados en información útil, rápida y comprensible para apoyar la respuesta ante incendios forestales.
